% \chapter{结论与展望}
\chapter{Conclusion and Outlook}

% 本文围绕 CMS 实验中 \(J/\psi + J/\psi + \phi\) 联合产生分析的论文写作需求，
% 搭建了按照实验分析流程组织的毕业论文结构。
% 该结构以绪论、物理动机、实验与数据、重建选择、效率与拟合、结果讨论和结论展望为主线，
% 便于后续将真实分析输出逐步填入。
% TO_APPROVE

This thesis structure has been developed to support the writing of an analysis of associated \(J/\psi + J/\psi + \phi\) production in the CMS experiment, organized according to the workflow of an experimental particle-physics analysis.
The structure follows the sequence of introduction, physics motivation, experiment and data samples, reconstruction and selection, efficiencies and fitting, results and discussion, and conclusion and outlook, allowing genuine analysis outputs to be incorporated progressively in later stages.

% 后续工作重点包括：
% TO_APPROVE

The main priorities for subsequent work include:

\begin{itemize}
  % \item 补全真实数据样本、触发和 MC 信息；
  \item completing the information on the real data sample, trigger, and MC samples;
  % \item 固化最终候选选择与信号提取方案；
  \item finalizing the candidate selection and signal extraction strategy;
  % \item 填入最终效率、系统误差与结果表格；
  \item filling in the final efficiency, systematic-uncertainty, and result tables;
  % \item 将外部工作区生成的图表按章节整理并补充来源说明。
  \item organizing figures produced in external workspaces by chapter and adding the corresponding provenance information.
\end{itemize}

% 在最终版论文中，本章还应回到物理问题本身，概括本分析对 \(J/\psi + J/\psi + \phi\)
% 联合产生研究能够提供的实验信息及其局限。
% TO_APPROVE

In the final version of the thesis, this chapter should return to the underlying physics question and summarize the experimental information that this analysis can provide on associated \(J/\psi + J/\psi + \phi\) production, together with its limitations.
